% Source: https://artofproblemsolving.com/community/c3h2658431_kappek_kmo2
% Stable identity: aops-topic-2658431-q01
Let the prime be $q$, so $x\ne780$ and $x^2=q(780-x)$. Then $q\mid x^2$, so $q\mid x$; write $x=qk$ with $k$ an integer. Dividing $q^2k^2=q(780-qk)$ by $q$,
\[qk(k+1)=780=2^2\cdot3\cdot5\cdot13.\]
So $q\in\{2,3,5,13\}$ and $k(k+1)=\dfrac{780}{q}\in\{390,260,156,60\}$. A product of two consecutive integers:
\begin{itemize}
\item $390$: $19\cdot20=380<390<420=20\cdot21$, impossible;
\item $260$: $15\cdot16=240<260<272=16\cdot17$, impossible;
\item $156=12\cdot13=(-13)(-12)$, so $k=12$ or $k=-13$;
\item $60$: $7\cdot8=56<60<72=8\cdot9$, impossible.
\end{itemize}
So $q=5$ and $x=5k\in\{60,-65\}$. Check: $\dfrac{60^2}{720}=5$ and $\dfrac{65^2}{845}=\dfrac{4225}{845}=5$. The answer is $\boxed{x=60\text{ or }x=-65}$.
